\chapter{Discussion}
\section{A feature decision depends on the fitted model}
The completed experiment began with an apparently simple reason to reduce the feature set: radius of gyration and SASA are strongly correlated. The comparison shows why that observation alone is insufficient. Correlation summarizes a relationship over the whole sample, whereas clustering assigns individual observations using the geometry of all retained features. The assignments can change when a feature is removed even if that feature shares much of its variation with another one.

The three-component mixtures retained high agreement in this pilot, both across feature sets and across the tested initializations. The other component counts did not give the same pattern. This does not select three as the physical number of conformational states. It shows that the answer to the feature-removal question depends on how the data are represented by the chosen model. Keeping all tested component counts makes this dependence visible instead of reporting only the most stable case.

The fixed fitting sample helps isolate the comparison but also limits it. Both feature sets use the same observations and the same scaling for shared features. Differences are therefore not caused by selecting a new sample for each fit. However, the experiment says nothing directly about how those results would change with another fitting subset. Initializations and sample selection are different sources of uncertainty, and only the former was varied here.

\section{Connecting the extension to the published workflow}
The primary study uses density-peak discovery, consensus and reweighting to characterize conformational families \autocite{cadenelli2026}. The GMM pilot examines a narrower statistical question using the available descriptors. Its agreement scores neither reproduce the paper's family populations nor test the reported difference between Aβ40 and Aβ42. They also give no direct information about aggregation rates or toxicity.

A common difficulty links the two approaches: a numerical family still needs a physical interpretation. A stable Gaussian partition might summarize a broad descriptor distribution rather than separate metastable states. The consensus baseline likewise compares marginal descriptor distributions and can miss differences in their joint dependence. Structures, contacts and alternative representations would help examine these possibilities. An independently constrained ensemble could provide another comparison, provided the structures, weights and observable definitions are compatible \autocite{borthakur2025}.

\section{Sampling, weighting and the next comparison}
The temporal diagnostics limit the uncertainty claims available from the present data. Randomly drawing fitting frames does not remove correlation from the process that generated them. All replica--descriptor series fail at least one estimator check, so the small provisional ESS values should be treated as warnings rather than calibrated effective counts. Repeating fits across fitting subsets and time blocks would address sampling sensitivity more directly \autocite{grossfield2019}.

Weighting introduces a further issue. Grid completeness establishes that the FES can be queried, but not that its scale or construction is appropriate for the chosen estimator. Once those details are verified, unequal weights may concentrate the estimate on relatively few observations. That effect must be assessed separately from seed variation and temporal dependence. With only three replicas, uncertainty based on differences between replicas will also need cautious interpretation.

Saved labels, source tables and numerical tests allow the completed calculations to be checked. Equilibrium family populations remain dependent on the missing FES provenance. This boundary keeps further use of the project tied to its evidence \autocite{plumed2019,barker2022}.
